\documentclass[../../main.tex]{subfiles}
\begin{document}

\begin{center}
  {\fontsize{48pt}{58pt}\selectfont \textbf{2001}\normalsize\textbf{年度}}
\end{center}

\begin{table}[htb]
  \centering
  \caption{出題テーマと難易度}
  \begin{tabular}{|c|p{9cm}|c|}
    \hline
    \textbf{問題番号} & \textbf{テーマ} & \textbf{難易度} \\
    \hline
    第1問 & 四面体の外接球の半径 & \\
    \hline
    第2問 & 積分方程式を満たす関数の一意性 & \\
    \hline
    第3問 & 双曲線と線分が囲む面積の比の単調性 & \\
    \hline
    第4問 & 複素数の漸化式と円周上の点 & \\
    \hline
    第5問 & ビーカーの水を合わせる操作 & \\
    \hline
    第6問 & コイン投げで動く二点の到達位置 & \\
    \hline
  \end{tabular}
\end{table}

\begin{table}[htb]
  \centering
  \caption{解答}
  \renewcommand{\arraystretch}{1.4}
  \begin{tabular}{|c|c|p{9cm}|}
    \hline
    \textbf{問題番号} & \textbf{小問} & \textbf{答え} \\
    \hline
    第1問 & -- & $r=\dfrac{\sqrt{13}}{3}$ \\
    \hline
    第2問 & -- & 条件は$b-2\ne\pm a$．そのとき
    $f(x)=\dfrac{2}{2-(a+b)}(\sin x+\cos x)$ \\
    \hline
    第3問 & -- & $t>1$で$c'(t)<0$（証明問題） \\
    \hline
    第4問 & (1) & 円$C$の中心は$\dfrac{1}{2}$，
    半径は$\dfrac{\sqrt{5}}{2}$ \\
    \cline{2-3}
    & (2) & すべての$b_n$が円$C$上にある（証明問題） \\
    \hline
    第5問 & (1) & $x<\dfrac{1}{3}$なら元の$x$のビーカーは途中で除かれるか，
    残って水量が増える（証明問題） \\
    \cline{2-3}
    & (2) & $x>\dfrac{2}{5}$なら元の$x$のビーカーは水量$x$のまま残る
    （証明問題） \\
    \hline
    第6問 & (1) & $X_{n+1}=2^n-X_n$ \\
    \cline{2-3}
    & (2) & $X_n=\dfrac{1}{3}\left\{2^n+2(-1)^n\right\}$ \\
    \cline{2-3}
    & (3) & $a$の平均は$\dfrac{2n}{3}-\dfrac{1}{9}
    \left\{1-\left(-\dfrac{1}{2}\right)^n\right\}$ \\
    \hline
  \end{tabular}
\end{table}

\paragraph{試験形態}
6問．試験時間は要確認．

\end{document}
